\documentclass[12pt]{article}
\usepackage[utf8]{inputenc}
\usepackage{amsmath, amssymb} % Math packages
\usepackage{geometry}
\usepackage{hyperref} % Clickable links
\usepackage{graphicx} % Images

\geometry{a4paper, margin=1in}

\title{Ottimizzazione combinatoria}
\author{Matteo}
\date{\today}

\begin{document}

\maketitle

\section{Introduct}

n questo corso dimostreremo che per sviluppare una soluzione per il supporto alle deci-
sioni sono necessari i seguenti passi:
\begin{itemize}
\item modellare matematicamente il problema che si vuole risolvere;
\item identificare la sua complessità;
\item individuare o progettare l’algoritmo più adatto, usando le tecniche di soluzione in
relazione al modello e alla complessità del problema.
\item Mostreremo come problemi molto semplici da formulare possono risultare molto
difficili da risolvere e come problemi molto complessi da descrivere possono essere
in realt`a molto facili da risolvere.
\item Partiremo definendo cosa `e un problema, un modello e un algoritmo.

\end{itemize}

\textbf{ricerca operativa}: Disciplina della matematica applicata che fornisce metodi scien-
tifici quantitativi per la formalizzazione e risoluzione di problemi di ottimizzazione che si
presentano in strutture organizzate complesse

molte discipline diverse all'interno della ricerca operativa(informatica, matematica, economia ecc...)

due modi di approcciarsi ai problemi:

\textbf{approccio qualitativo}: si fanno scelte in base alle esperienze passate e in modo democratico. scelte basate anche su criteri soggettivi e concetti non matemattici come affidabilita e sicurezza.

\textbf{approccio quantitativo}: si analizzano tutte le scelte(usando algoritmi). scelte basate solo su numeri e statistiche.

si puo ragionare sui problemi di ottimizzazione sia a livello strategico che a livello operativo

\textbf{livello strategico}: definizione e valutazione delle alternativi in chiave probabilistica

\textbf{livello operativo}: definizione della prassi operativa in base alle scelte strategiche fatte. si operaa in regime deterministico.

difficoltà di un problema= tempo necessario (a un algoritmo) a risolverlo all'aumentare della sua dimensione.

due tipi di problemi:
\begin{itemize}
\item problemi polinomiali(P): il tempo dell'algoritmo cresce cosi: $0(1),0(n),0(n^2).....$
\item problemi non polinomiali(NP): il tempo cresce in maniera esponenziale:$0(2^n)$
\end{itemize}

La Ricerca Operativa fornisce metodi per
rappresentare (modelli) e risolvere (algoritmi)
problemi di ottimizzazione.

\textbf{modello}: rappresentazione semplificata di un sistema reale in grado di dare output(prestazioni) rispetto a certi input(decisioni)

esistono modelli fisici e modelli matematici

step per formulare un problemaa:
\begin{enumerate}
\item definizione del problema
\item definizione delle variabili
\item definizione della funzione obbiettivo come combinazione delle variabili decisionali
\item definizione dei vincoli come combinazione delle variabili decisionali.
\end{enumerate}

una volta fatto il modello, lo si verifica e calibra e vi si esperimenta su.

poi si generano soluzione mediante un algoritmo.

poi sii presentano i risultati e si fanno opportune verifiche sulle ipotesi sul modello e obbiettivi e soluzioni.

esistono anche problemi non necessariamente matematici(costi, logistica, pianificazione, gesstione....)

\section{problemi di ottimizzazione}

Un problema (di ottimizzazione) non è nient’altro che una domanda,
espressa in termini generali, ma la cui risposta dipende da un certo
numero di parametri e variabili.

Un’istanza del problema P si ottiene specificando dei valori concreti per
tutti i parametri del problema (ma non per le variabili!).




\end{document}
